% 合成 vs ShezhenV3 全量对比
\begin{table}[t]
\centering
\caption{从合成阶段 2 到 ShezhenV3 全量验证的演进。}
\label{tab:synthetic-vs-tcm-zh}
\small
\resizebox{\columnwidth}{!}{%
\begin{tabular}{lcc}
\hline
量 & 合成（阶段 2） & ShezhenV3（全量） \\
\hline
标定 $n$ & 240 & 1{,}144 \\
清晰/模糊中位数 & 0.819 / 0.191 & 0.516 / 0.414（ROI） \\
$\tau$ & 0.505 & 0.465 \\
M5 Spearman $\rho$ & $\approx$0.75 & 0.363 \\
每 seed 测试帧 & 500（仿真 $Q$） & 500（真实分数） \\
B2 M1 p50 (ms) & 364.7 & 208.3 \\
B2 M2 有效帧率 & 0.988 & 0.604 \\
\hline
\end{tabular}%
}
\end{table}

\paragraph{迁移解读。}
真实纹理使 $Q_{\mathrm{img}}$ 绝对值下降；Spearman 降低是因为清晰/模糊组重叠增大（见 \S\ref{sec:tcm-iqa-zh}）。
B2 仍保持并略强化有效帧优势，且相对 B0 改善了 RTT 中位数，说明路由架构可超出合成直方图泛化。
ROI 与学习型 B3 已在正文中实现（\S\ref{sec:roi-zh}、\S\ref{sec:b3-learned-zh}）；后续可探索 ROI 与混合 flag 规则的组合。

\paragraph{JSONL 审计样例。}
补充材料 Listing~1 为 Franka 辅轨的脱敏 JSONL 行，字段与 \S\ref{sec:closed-loop-zh} 一致：
\begin{verbatim}
{"trial_id":1,"q_img":0.8075,"flags":[],
 "retry_count":0,
 "route_decision":"upload_cloud",
 "latency_ms":15.89}
\end{verbatim}
完整日志另含 ISO~8601 时间戳，由 \texttt{validate\_jsonl.py} 校验。

\paragraph{算力预算。}
全验证打分（1{,}144 对）加三 seed 主矩阵，在笔记本 CPU 上可在十分钟内完成，便于新 TCM 批次到达后频繁重标定。
